\clearpage
\SourceLeaf{264}
\begin{center}
{\large SECTIO QUARTA.}\par\medskip
\pdfbookmark[1]{De determinatione orbitarum, habita ratione perturbationum.}{liber-ii-sectio-iv}
\textit{De determinatione orbitarum, habita ratione perturbationum.}
\end{center}
\Needspace{6\baselineskip}
\begin{center}\hypertarget{native-264-article190}{}\pdfbookmark[2]{190.}{article-190}190.\end{center}
Perturbationes, quas planetarum motus per actionem planetarum reliquorum
patiuntur, tam exiguae lentaeque sunt, ut post longius demum temporis
intervallum sensibiles fiant: intra tempus brevius --- vel adeo, prout circumstantiae
sunt, per revolutionem integram unam pluresve --- motus tam parum
differet a motu in ellipsi perfecta secundum leges \textsc{Kepler}i exacte descripta,
ut observationes deviationem indicare non valeant. Quamdiu res ita se habet,
operae haud pretium esset, calculum praematurum perturbationum suscipere,
sed potius sufficiet, sectionem conicam quasi osculatricem observationibus adaptare:
dein vero, postquam planeta per tempus longius accurate observatus
est, effectus perturbationum tandem ita se manifestabit, ut non amplius possibile
sit, omnes observationes per motum pure ellipticum exacte conciliare;
tunc itaque harmonia completa et stabilis parari non poterit, nisi perturbationes
cum motu elliptico rite iungantur.

Quum determinatio elementorum ellipticorum, cum quibus perturbationes
iungendae sunt, ut observationes exacte repraesententur, illarum cognitionem
supponat, vicissim vero theoria perturbationum accurate stabiliri nequeat, nisi
elementa iam proxime cognita sint: natura rei non permittit, arduum hoc
negotium primo statim conatu perfectissime absolvere, sed potius perturbationes
et elementa per correctiones alternis demum vicibus pluries repetitas ad summum
praecisionis fastigium evehi poterunt. Prima itaque perturbationum
\SourceLeaf{265}theoria superstruetur elementis pure ellipticis, quae observationibus proxime
adaptata fuerant: dein orbita nova investigabitur, quae cum his perturbationibus
iuncta observationibus quam proxime satisfaciat. Quae si a priori considerabiliter
discrepat, iterata perturbationum evolutio ipsi superstruenda erit,
quae correctiones alternis vicibus toties repetentur, donec observationes, elementa
et perturbationes quam arctissime consentiant.
